\documentclass[a4paper,11pt]{article}
\usepackage[margin=2.5cm]{geometry}
\usepackage{amsmath,amssymb}
\usepackage{booktabs}
\usepackage{hyperref}
\usepackage{xcolor}
\usepackage{listings}
\usepackage{enumitem}

\hypersetup{colorlinks=true, linkcolor=blue!50!black, urlcolor=blue!60!black}

\lstset{
  basicstyle=\ttfamily\small,
  breaklines=true,
  frame=single,
  framesep=4pt,
  rulecolor=\color{gray!50},
  backgroundcolor=\color{gray!5},
  columns=fullflexible,
}

\title{Parity Blindness, Phase 3, and the \texttt{repS} Concern \\
       \large How a residual faithfulness defect in the symcoder encoder
               was diagnosed and closed by switching from \texttt{repS}
               to \texttt{rep}}
\author{Echinocoder / symcoder working notes}
\date{2026-05-25 (initial), 2026-05-26 (resolution recorded)}

\begin{document}
\maketitle

\begin{abstract}
This document captures an investigation into a class of faithfulness
defects in the symcoder encoder.  Phase 1 + Phase 2 of the existing
encoder were found to be parity-blind on 31 specific
``shadow'' events.  A new ``Phase 3'' was added in which the
deterministically faithful simplicial-complex multiset embedder is
applied to the alignment table $T$ of the event.  In its first
implementation Phase 3 closed 29 of the 31 cases.  The two remaining
shadows turned out to be a symptom of the more interesting structural
problem the title alludes to: the alignment table was being built over
\texttt{repS}, the pre-pruned ``short'' atom set, whose pruning had
been justified \emph{only} under the (now-known-false) assumption that
pair-row marginals are sufficient.  Switching Phase 3 to use the
unpruned \texttt{rep} closed all 31 cases.  The document records the
chain of reasoning, the implementation, and the residual open
questions, so that the work can be picked up cold from the repository.
\end{abstract}

\tableofcontents

\section{What problem are we solving?}

\subsection{The encoder's contract}
Echinocoder's symcoder encodes a labelled event $E$ (a tuple of vectors
in $\mathbb{R}^d$, labelled by particle names) as a fixed-length real
vector that is invariant under:
\begin{itemize}[nosep]
  \item the \emph{label-permutation group} $H = S_{n_1} \times S_{n_2}
        \times \cdots$ acting on particle labels of each type;
  \item the rotation group $SO(d)$ acting on the underlying vectors;
  \item (for $d=3$ specifically) \emph{not} the full $O(3)$: parity
        (reflection) is supposed to be preserved by the encoder because
        the operator set includes the pseudoscalar $\varepsilon_3$.
\end{itemize}

\noindent
End users use only the encoder; decoders exist solely as validation
tools that confirm round-trips lose nothing other than the intended
$H \times SO(d)$ ambiguity.

\subsection{The defect: parity blindness on structured events}
For ``generic'' real-valued events the existing two-phase encoder
(Phase~1: per-orbit value multisets; Phase~2: per-row-pair joint value
multisets) is empirically parity-faithful: $\text{encode}(E)$ differs
from $\text{encode}(P(E))$ whenever $E$ and $P(E) = \{-v : v \in E\}$
lie in distinct $SO(3) \times H$ orbits.

However, for certain \emph{structured} events --- typically with
antipodal muons (e.g.\ $q = -p$) and integer-valued electron
coordinates --- the encoder produces \emph{exactly the same output}
for $E$ and $P(E)$, despite the two being in different $SO(3) \times H$
orbits.  This was discovered by the empirical search in
\texttt{symcoder/experiments/encoding\_fidelity\_search.py} and pinned
by 31 regression test cases in
\texttt{symcoder/tests/test\_parity\_blindness\_shadows.py},
parametrised over the fixture
\texttt{symcoder/tests/\_parity\_shadow\_events.py}.

The 31 ``shadow events'' are events for which
$\text{encode}(E) = \text{encode}(P(E))$ exactly, yet the Gram-matrix
plus signed-determinant invariant (Weyl's First Fundamental Theorem for
$SO(n)$) distinguishes them.  All 31 are chiral in the strict sense:
the chirality precondition in the test file confirms this case-by-case.

\subsection{Why Phase 1 + Phase 2 alone is insufficient}
For any antisymmetric operation (such as $\varepsilon_3$), the
multiset of orbit values is intrinsically palindromic under negation
(every value $v$ appears together with $-v$ in the orbit).  Hence
Phase~1 \emph{never} distinguishes parity for these rows.

Phase~2 (per-row-pair joint value multisets) usually breaks this
palindrome on generic inputs but does not always: for the 31 shadow
events the joint multisets are palindromic in the antisymmetric
coordinate too.

This is a property of \emph{pair-row} marginals.  The information
needed to distinguish chirality lives at higher arity (joint structure
across $\ge 3$ rows simultaneously); pair marginals miss it on
structured inputs.

\section{The attempted fix: Phase 3 simplicial-complex embedding}

\subsection{Architecture}
A new encoder factory has been wired into
\texttt{symcoder/encode.py}:

\begin{lstlisting}
encode(plan, event,
       orbit_factory=...,        # Phase 1
       phase2_factory=...,       # Phase 2
       phase3_factory=Phase3SimplicialEncoderFactory()  # default on
)
\end{lstlisting}

The new factory builds a \texttt{Phase3SimplicialEncoder} (defined in
\texttt{symcoder/encoders/phase3\_simplicial.py}) which:

\begin{enumerate}[nosep]
  \item Computes the full \emph{alignment table}
        $T \in \mathbb{R}^{|G| \times |\text{repS}|}$ by direct
        evaluation: $T_{g,i} = \text{eval}(g \cdot \text{atom}_i, E)$
        where the rows are indexed by all group elements
        $g \in G$ and the columns by all canonical \texttt{repS} atoms.
  \item Treats $T$ as a \emph{multiset of $|G|$ tuples of dimension
        $|\text{repS}|$} (i.e.\ the column ordering is forgotten).
  \item Feeds this multiset to the existing top-level
        \texttt{C0HomDeg1\_simplicialComplex\_embedder\_2} simplicial
        embedder, which produces a deterministically faithful
        embedding of any multiset of $n$ $k$-tuples in
        $2nk + 1 - k$ reals.
  \item Concatenates that vector to the end of the existing
        Phase~1 + Phase~2 output, with a new \texttt{Phase3Tree}
        segment-descriptor entry so users can slice it off if they
        only want the smooth Phase~1+2 prefix.
\end{enumerate}

\subsection{First-pass result: 29 out of 31 shadow events fixed}
With \texttt{phase3\_factory} on by default and the alignment table
built over \texttt{repS}, 29 of the 31 parity regression tests pass:
$\text{encode}(E) \neq \text{encode}(P(E))$ at the required tolerance.
A companion test
(\texttt{test\_parity\_flip\_no\_phase3\_still\_blind}) explicitly
passes \texttt{phase3\_factory=None} and pins all 31 of the original
failures as \texttt{xfail} so that the historical Phase-1+2-only bug
remains visible.

The two surviving failures (shadows 16 and 24) led to the deeper
finding described in \S4.

\section{The two intermediate failures: shadows 16 and 24}

\emph{Note (post-resolution):} the failures discussed in this section
were observed when Phase 3 used \texttt{repS}.  After switching Phase 3
to use the unpruned \texttt{rep} (see \S6), both shadows now pass.
The diagnostic content below is retained because it explains
\emph{why} the \texttt{repS} pruning was wrong --- which is also the
reason \texttt{repS} cannot in general be trusted as the column set of
the alignment table.


\subsection{What the diagnostic shows}
For these two shadow events, the alignment tables $T(E)$ and
$T(P(E))$ are \emph{the same multiset of 12 column vectors},
bit-for-bit identical (zero difference at IEEE precision after lex
sorting).  Hence \emph{any} encoder whose input is the multiset of
columns of $T$ must produce identical output for these two events.

This is not a numerical artefact of the simplicial encoder's
piecewise-linear corner structure.  It was checked with
\texttt{np.array\_equal} on lex-sorted rows.

\subsection{Why this is mathematically possible}
The label group $H$ acts on $T$ (viewed as a function
$G \to \mathbb{R}^{|\text{repS}|}$) by \emph{left translation}:
$T(\sigma \cdot E)(g) = T(E)(\sigma^{-1} \cdot g)$.  An $H$-invariant
encoding need only be invariant under this left-translation action.

The multiset-of-columns is invariant under \emph{any} permutation of
columns --- including permutations that are \emph{not} left
translations by any $\sigma \in H$.  So
``multiset of columns'' is a \emph{strictly coarser} equivalence than
``function modulo left translation by $H$''.

In other words: two distinct functions $T_1, T_2 : G \to
\mathbb{R}^{|\text{repS}|}$ can give the same column-multiset without
being in the same $H$-orbit.  For shadows 16 and 24 this happens; the
encoder consequently cannot distinguish them.

\subsection{The geometric cause}
For these events $T(E)$ has a high-symmetry column pattern: many
columns share the same Phase-1-style $(\text{mag}, \text{dot})$
profile, and within each such cluster the $\varepsilon_3$ values come
in $\pm$ pairs.  A column-permutation that pairs each $+$-eps column
with its $-$-eps counterpart, while leaving the $(\text{mag}, \text{dot})$
profile per cluster unchanged, takes $T(E)$ to $T(P(E))$ as a
multiset --- but this permutation is \emph{not} of the form ``left
translation by some $\sigma \in H$.''

\section{The \texttt{repS} concern (the key insight at the end of session)}

\subsection{Historical pruning}
The \texttt{repS} set used in the current code is the result of a
historical pruning process:

\begin{itemize}[nosep]
  \item \texttt{rep} (original): \emph{all} mags, all (dissimilar)
        dot pairs, and all (dissimilar) $\varepsilon_3$ triples
        across the labels of the plan.  By Weyl's First Fundamental
        Theorem this is information-complete for $SO(3)$ invariants
        of the labelled vector tuple.
  \item \texttt{repL} (intermediate): a longer-than-\texttt{repS}
        intermediate pruning, no longer present in the codebase.
  \item \texttt{repS} (current, ``short''): further pruned by
        observing that many pairs of \texttt{rep} rows are
        $G$-orbit-equivalent (e.g.\ $(\text{dot}(a,b), \text{dot}(c,d))$
        and $(\text{dot}(a,d), \text{dot}(b,c))$ both have the same
        $G$-orbit at the pair level), so only one of each duplicate
        pair needs to be stored.
\end{itemize}

\subsection{Why the pruning is now suspect}
The pruning from \texttt{rep} $\to$ \texttt{repL} $\to$ \texttt{repS}
was justified \emph{under the assumption that Phase 1 + Phase 2
multisets are sufficient}: the duplicates being eliminated were
duplicates \emph{at the pair-row marginal level}.

We now know (from the 31 shadow events) that pair-row marginals are
\emph{not} sufficient for faithful encoding.  Therefore the pruning
that yielded \texttt{repS} was based on a false premise: rows that
look ``duplicate'' at the pair-marginal level may carry distinct
joint information once we look at higher-arity marginals (such as the
full alignment table).

In particular, feeding \texttt{repS}-derived $T$ to the simplicial
encoder makes the same mistake the pair-only pruning made: it throws
away rows whose presence might break the column-multiset palindrome.

\subsection{Concrete construction of \texttt{rep}}
The full \texttt{rep} can be regenerated from a plan as the
concatenation of:

\begin{enumerate}[nosep]
  \item $[\text{mag}(v) : v \in \text{labels}]$ --- one mag per label.
  \item $[\text{dot}(a,b) : \{a,b\} \subset \text{labels},\ a \neq b]$
        --- one dot per \emph{unordered} pair of distinct labels.
        Iterating over \emph{combinations} (not permutations) avoids
        double-counting; the dot constructor already
        canonicalises its arguments.
  \item $[\varepsilon_3(a,b,c) : \{a,b,c\} \subset \text{labels},\
        a,b,c \text{ distinct}]$ --- one $\varepsilon_3$ per
        \emph{unordered} triple of distinct labels.  Again iteration
        over combinations is enough; the $\varepsilon_3$ constructor
        is antisymmetric and canonicalises.
\end{enumerate}

For a 3-electron, 2-muon plan:
\[
  |\text{rep}| = \underbrace{5}_{\text{mags}}
               + \underbrace{\binom{5}{2}}_{\text{dots} = 10}
               + \underbrace{\binom{5}{3}}_{\varepsilon_3 = 10}
               = 25,
\]
versus $|\text{repS}| = 8$ currently used.  $T$ becomes
$|G| \times |\text{rep}| = 12 \times 25 = 300$ entries.  The
simplicial embedder size then becomes $2 \cdot 12 \cdot 25 + 1 - 25 =
576$ reals, replacing the current 185.

\subsection{The hypothesis (now confirmed)}
Using \texttt{rep} (not \texttt{repS}) in the Phase 3 alignment table
\emph{should} make the column-multiset of $T$ injective on
$SO(3) \times H$ orbits for the two stubborn shadow events.  When
this swap was implemented, all 31 shadow events --- including 16 and
24 --- did indeed become distinguishable by the encoder.  See \S6 for
the implementation and the resulting test status.

\medskip
\noindent
\textbf{An important caveat that has not been disproven.}  Even with
\texttt{rep}, the column-multiset of $T$ is a strictly coarser
equivalence than ``function modulo left-translation by $H$''.  We have
no theorem that says \texttt{rep}'s column-multiset is a complete
invariant for $SO(3) \times H$ in general.  All we know is that on the
31 events of the current fixture it suffices; an adversarial event
constructed against the column-multiset of \texttt{rep} is still
mathematically possible (just harder to find than against
\texttt{repS}).  See \S\ref{sec:open-questions}.

\section{The resolution: switching Phase 3 from \texttt{repS} to \texttt{rep}}
\label{sec:resolution}

\subsection{What was changed in the code}

\begin{enumerate}[nosep]
  \item A new function \texttt{rep(context, operations)} was added to
        \texttt{symatom/rep.py}, exported from \texttt{symatom/\_\_init\_\_.py}.
        It flattens \texttt{repS}'s FlavouredOperators into a list of
        individual \texttt{Atom}s: every atom across every G-orbit, with
        one atom per sign-canonical form.  For symmetric ops this is one
        atom per unordered label combination of the operation's rank; for
        antisymmetric ops it is one atom per unordered label combination
        (the $\pm$ pair is collapsed by storing only the
        \texttt{sign=+1} representative).
  \item \texttt{symcoder/encoders/phase3\_simplicial.py} was switched
        from \texttt{repS} to \texttt{rep}.  The alignment table is now
        $(|G|, |\text{rep}|)$ instead of $(|G|, |\text{repS}|)$.
  \item No other production code was touched.  Phase 1 and Phase 2
        continue to use \texttt{repS} internally, which is correct:
        the \texttt{repS} pruning was justified at the pair-marginal
        level, and Phase 1 + Phase 2 only consume pair marginals.  Only
        Phase 3, which claims full faithfulness, needed \texttt{rep}.
\end{enumerate}

\subsection{What this costs}

For the 3-electron, 2-muon parity-fixture plan:

\medskip
\noindent
\begin{tabular}{lrr}
\toprule
                                    & \texttt{repS}   & \texttt{rep} \\
\midrule
\#\,atoms ($=k$)                    & 8               & 25 \\
alignment table shape               & $12 \times 8$   & $12 \times 25$ \\
Phase 3 output size $2nk+1-k$       & 185 reals       & 576 reals \\
total encoder output (Ph1+Ph2+Ph3)  & $\approx 800$ reals & $\approx 1200$ reals \\
\bottomrule
\end{tabular}
\medskip

\noindent
The growth scales as roughly
$\binom{|\text{labels}|}{\text{rank of $\varepsilon$}}$ for the
$\varepsilon$ row count, which is the dominant cost for plans with
many particles and rank-3 (or higher) antisymmetric operations.  This
is a known and accepted cost; it is the price of taking the unpruned
information-complete set of rows.

\subsection{Test status after the change}

\begin{itemize}[nosep]
  \item \texttt{test\_parity\_flip\_distinguishable} (Phase 3 on,
        \texttt{rep}-based): \textbf{31 of 31 pass}.
  \item \texttt{test\_parity\_flip\_no\_phase3\_still\_blind}
        (Phase 3 off, pinning the historical bug): 31 of 31 xfail
        (as expected).
  \item Total symcoder/symatom suite: 622 passed, 32 xfailed.
        The 32 xfails are 1 original alignment-decoder xfail + the
        31 historical-no-Phase-3 xfails.
\end{itemize}

\section{File locations and how to reproduce}

\subsection{Code under investigation}
\begin{itemize}[nosep]
  \item \texttt{symcoder/encode.py} --- top-level encode functions; Phase 3
        wired in with default-on \texttt{phase3\_factory} sentinel.
  \item \texttt{symcoder/encoders/phase3\_simplicial.py} --- the new
        \texttt{Phase3SimplicialEncoder} + factory.  Now uses the
        unpruned \texttt{rep}.
  \item \texttt{symatom/rep.py} --- defines both \texttt{repS} (the
        historical pruned set, still used by Phase 1 + Phase 2) and the
        new \texttt{rep} function used by Phase 3.
  \item \texttt{symcoder/describe.py} --- \texttt{SegmentInfo},
        \texttt{Phase3Tree}, \texttt{EncodingTree.phase3}.
  \item \texttt{C0HomDeg1\_simplicialComplex\_embedder\_2\_for\_array\_of\_reals\_as\_multiset.py}
        --- the underlying simplicial-complex multiset embedder
        (top-level Echinocoder file, not part of any package).
        Output size $2nk+1-k$ for input shape $(n,k)$.
\end{itemize}

\subsection{Tests}
\begin{itemize}[nosep]
  \item \texttt{symcoder/tests/test\_parity\_blindness\_shadows.py}:
    \begin{itemize}[nosep]
      \item \texttt{test\_parity\_flip\_distinguishable} (Phase 3 on):
            31/31 pass after the \texttt{repS} $\to$ \texttt{rep} fix.
      \item \texttt{test\_parity\_flip\_no\_phase3\_still\_blind}
            (Phase 3 off, xfail-pinned): 31/31 xfail.
    \end{itemize}
  \item \texttt{symcoder/tests/\_parity\_shadow\_events.py} ---
        auto-generated fixture; 31 shadow events as Python literals.
\end{itemize}

\subsection{Experimental discovery script}
\texttt{symcoder/experiments/encoding\_fidelity\_search.py} ---
brute-force search over event families; SO(3) and O(3) oracles based
on Gram-plus-signed-dets up to typed permutations of labels.
Re-runnable to discover further shadow events or to validate that a
fix really fixes \emph{all} known shadows.

\subsection{One-shot diagnostic}
To check whether $T(E)$ and $T(P(E))$ are the same multiset for any
event, see the snippet near the end of the chat: build the table with
\texttt{Phase3SimplicialEncoder.\_build\_table}, lex-sort the rows,
compare with \texttt{np.array\_equal}.

\section{Status and remaining open questions}

\subsection{Status today (after resolution)}
\begin{itemize}[nosep]
  \item Phase 3 is implemented, wired in, default-on, and uses the
        unpruned \texttt{rep} as the alignment-table column set.
  \item Test suite: \textbf{622 passed, 32 xfailed}.  The 32 xfails are
        1 pre-existing alignment-decoder xfail + 31
        historical-no-Phase-3 xfails.  All 31 parity-blindness
        regression cases pass with Phase 3 on.
  \item Both demos (\texttt{demo\_roundtrip.py},
        \texttt{demo\_shadow\_event.py}) still build; the
        parity-blindness demo PDFs are visibly distinguishable now
        that Phase 3 is on.
  \item No production code outside Phase 3 was changed.  Phase 1 and
        Phase 2 continue to use \texttt{repS}; that is correct because
        the \texttt{repS} pruning is valid at the pair-marginal level.
\end{itemize}

\subsection{Recommended follow-ups (no longer urgent, but worth doing)}

\begin{enumerate}
  \item \textbf{Stress-test the \texttt{rep} column-multiset
        faithfulness.}  The current fixture has 31 shadow events;
        \texttt{rep} clears all 31, but does not prove the
        column-multiset of $T(\text{rep})$ is a complete $SO(3)
        \times H$ invariant in general.  Extending
        \texttt{encoding\_fidelity\_search.py} to search against the
        \texttt{rep}-based encoder (rather than the \texttt{repS}
        one) might surface fresh adversarial events.
  \item \textbf{Document the dual role of \texttt{repS} vs \texttt{rep}.}
        The current docstrings explain why Phase 3 uses \texttt{rep},
        but the broader narrative ``\texttt{repS} for pair-marginal
        encoders only; \texttt{rep} for whole-table encoders that
        claim faithfulness'' should probably be surfaced in
        \texttt{SPEC.md} as a top-level design principle.
  \item \textbf{Per-row sparsity / sub-selection of \texttt{rep}.}
        \texttt{rep} grows as $\binom{|\text{labels}|}{\text{rank}}$
        which is painful for large plans.  Worth investigating
        whether there is a structured \emph{subset} of \texttt{rep}
        (call it \texttt{repM} for medium) which is still sufficient
        for whole-table faithfulness but smaller.  Any such subset
        must be characterised by a proof, not a heuristic, since the
        history of \texttt{repS} is a cautionary tale.
\end{enumerate}

\label{sec:open-questions}
\subsection{Open theoretical question}

\begin{itemize}[nosep]
  \item Multiset-of-columns of $T(\text{rep})$ is still a strictly
        coarser equivalence than ``function modulo left-translation
        by $H$''.  Even with \texttt{rep} fixed, palindromic
        column-multisets may remain possible in principle.
  \item If so, the next architectural lever is to encode $T$ as a
        function modulo $H$ (e.g.\ via canonical lex-min over the
        $|H|$ left-translations of $T$) rather than as a multiset.
        This is more corner-prone than simplicial but is the only
        obvious way to recover the lost group-labelling structure.
  \item Alternatively, augment $T$ with derived rows
        (e.g.\ row-products of antisymmetric and symmetric rows)
        whose presence breaks the column-multiset palindrome.
        Construction of provably non-palindromic augmentations is
        an open problem.
\end{itemize}

\label{sec:rep-construction}
\section*{Appendix: explicit \texttt{rep} construction recipe}

\noindent
\emph{Note:} this recipe was the design proposal during the
investigation.  The actual implementation in
\texttt{symatom/rep.py} achieves the same result by reusing existing
machinery: it constructs a \texttt{FlavouredOperator} for every
(operation, flavour) pair, then iterates over each
\texttt{FlavouredOperator.atoms\_one\_per\_sign()}.  Whether you read
the recipe as ``iterate combinations across all labels for each
operation rank'' (this appendix) or ``flatten every G-orbit in
\texttt{repS}'' (the implementation), the resulting atom list is
identical.

\bigskip
\noindent
Given a \texttt{Plan} with \texttt{context.types[i].labels} for each
particle type $i$, and any choice of operations (rank-1 symmetric,
rank-2 symmetric, rank-3 antisymmetric, etc.):

\begin{enumerate}[nosep]
  \item Let \texttt{labels} be the concatenated tuple of all labels
        across all types.
  \item Emit one $\text{mag}(\ell)$ atom for each $\ell \in
        \texttt{labels}$.
  \item Emit one $\text{dot}(\ell_i, \ell_j)$ atom for each unordered
        pair $\{\ell_i, \ell_j\} \subseteq \texttt{labels}$ with
        $\ell_i \neq \ell_j$
        (i.e.\ \texttt{itertools.combinations(labels, 2)}).
  \item Emit one $\varepsilon_3(\ell_i, \ell_j, \ell_k)$ atom for each
        unordered triple $\{\ell_i, \ell_j, \ell_k\} \subseteq
        \texttt{labels}$ with the three labels pairwise distinct
        (i.e.\ \texttt{itertools.combinations(labels, 3)}).
\end{enumerate}

\noindent
The dot and $\varepsilon_3$ atom constructors are responsible for
their own canonical-form sorting (lex order for symmetric ops, sign
extraction for antisymmetric ops), so iterating over
\texttt{combinations} (unordered) is sufficient.

For multi-type plans the labels of different types are concatenated;
no special treatment is required because \texttt{mag}, \texttt{dot},
and $\varepsilon_3$ operate on labels independent of which type each
label belongs to.

\end{document}
